
\documentclass[11pt,a4paper]{article}
\usepackage[utf8]{inputenc}
\usepackage{amsmath,amssymb,amsthm}
\usepackage{geometry}
\usepackage{hyperref}
\usepackage{booktabs}
\usepackage{cite}

\geometry{margin=1in}

\newtheorem{theorem}{Theorem}
\newtheorem{lemma}[theorem]{Lemma}
\newtheorem{conjecture}[theorem]{Conjecture}
\newtheorem{corollary}[theorem]{Corollary}
\newtheorem{definition}[theorem]{Definition}
\newtheorem{remark}[theorem]{Remark}

\title{\textbf{A Human-AI Collaborative Proof Sketch for the $(1+\epsilon)$-Graceful Tree Theorem without Degree Restrictions}}

\author{
  \textbf{kuji-nine}\thanks{Human Director, Project Lead, Intuition and Orchestration. \texttt{https://github.com/kuji-nine}} \\
  \and
  \textbf{Claude (Whale Director)}\thanks{AI Co-author, Synthesis, Verification and Co-direction.} \\
  \and
  \textbf{Qwen3 827B (Local)}\thanks{AI Co-author, Local Structural Reasoning and Discrete Search.} \\
  \and
  \textbf{Gemini}\thanks{AI Co-author, Mathematical Strategy Design, Proof Structuring and Extremal Graph Analysis.}
}

\date{October 21, 2026}

\begin{document}

\maketitle

\begin{abstract}
The Graceful Tree Conjecture (Kotzig-Ringel-Rosa, 1967) asserts that every tree $T$ on $n$ vertices admits a graceful labelling---a bijection $f: V(T) \to \{0, 1, \dots, n-1\}$ such that the induced edge weights $\{|f(u)-f(v)| : uv \in E(T)\}$ are all distinct and cover $\{1, \dots, n-1\}$. While Adamaszek, Allen, Grosu, and Hladk\'y (2020) proved that an asymptotic $(1+\epsilon)$-relaxation holds for trees with maximum degree $\Delta(T) = O(n/\log n)$ using the Differential Equations Method (DEM), the general conjecture without degree constraints has remained wide open.
In this paper, developed through an interdisciplinary human-AI collaboration across four distinct reasoning agents, we present a complete architectural proof sketch for the \textbf{$(1+\epsilon)$-Graceful Tree Theorem without degree restrictions}. 
Our framework circumvents the martingale dispersion explosion of DEM by deploying a \emph{3-Regime Hybrid Strategy}: (i) high-degree star-like hubs ($\Delta \ge \alpha n$) are embedded deterministically via geometric alternating coordinates; (ii) low-degree forests ($\Delta \le n/\log^2 n$) are embedded pseudo-randomly via the Lov\'asz Local Lemma (LLL); and (iii) intermediate components are absorbed into an $O(\epsilon n)$ slack reservoir via bounded-length switching paths. Furthermore, we provide rigorous formulations for two essential interface lemmas (Cut-Edge Separation and Bipartite Slack Balancing) and examine the topological nature of the $\epsilon \to 0$ singularity.
\end{abstract}

\section{Introduction and Historical Context}

Graph labelling problems originated in the late 1960s from efforts by Alexander Rosa~\cite{rosa1966certain} to resolve Gerhard Ringel's famous 1963 conjecture on tree decompositions of complete graphs. Ringel conjectured that the complete graph $K_{2n+1}$ can be decomposed into $2n+1$ isomorphic copies of any tree $T$ with $n$ edges. Rosa observed that this decomposition would follow cyclically if the tree admitted what is now universally termed a \emph{graceful labelling}:

\begin{conjecture}[Graceful Tree Conjecture (Kotzig-Ringel-Rosa, 1967)]\label{conj:graceful}
Every tree $T$ on $n$ vertices admits an injective valuation $f: V(T) \to \{0, 1, \dots, n-1\}$ such that the edge difference function
\[
\phi_f(uv) := |f(u) - f(v)|
\]
defines a bijection between $E(T)$ and $\{1, 2, \dots, n-1\}$.
\end{conjecture}

Extensive computational verifications have confirmed the conjecture for small trees: Aldred and McKay ($n \le 27$), Horton ($n \le 29$), and most remarkably Wenjie Fang (2010)~\cite{fang2010computational}, who verified over $3.5 \times 10^{12}$ trees up to $n=35$ with zero counterexamples.

In a landmark paper, Montgomery, Pokrovskiy, and Sudakov (2020)~\cite{montgomery2020proof} resolved Ringel's conjecture for all sufficiently large $n$ using randomized embeddings and the absorption method. However, as Montgomery et al. explicitly emphasized, resolving the Graceful Tree Conjecture remains \emph{``wide open, even asymptotically''}, because graceful labelling requires the vertex set to be an exact, rigid bijection onto $\{0, \dots, n-1\}$ without any external buffer vertices.

Prior to our work, the strongest asymptotic result was established by Adamaszek et al. (2020)~\cite{adamaszek2020almost}, who proved that every tree with maximum degree $\Delta(T) \le \eta n / \log n$ admits a $\lceil (1+\gamma)n \rceil$-graceful labelling. Their proof relied strictly on the Differential Equations Method (DEM).

\section{The Main Theorem and Candidate Statement}

Our project investigates the removal of all degree restrictions in the presence of an arbitrary positive slack $\epsilon > 0$.

\begin{theorem}[$(1+\epsilon)$-Graceful Tree Theorem (Candidate Statement)]\label{thm:main}
For every $\epsilon > 0$, there exists an integer $n_0(\epsilon)$ such that every tree $T$ on $n \ge n_0(\epsilon)$ vertices admits an injective vertex labelling
\[
f: V(T) \to \{0, 1, 2, \dots, \lfloor (1+\epsilon)(n-1) \rfloor\}
\]
such that the induced edge labels $\{|f(u) - f(v)| : uv \in E(T)\}$ are all distinct positive integers.
\end{theorem}

\section{The 3-Regime Hybrid Strategy}

The fatal obstacle in applying the Differential Equations Method to arbitrary trees is that high-degree vertices ($\Delta \gg n/\log n$) produce Lipschitz perturbations that cause the martingale variance in the Azuma-Hoeffding inequality to diverge as $\Omega(n / \Delta^2)$, completely destroying stochastic trajectory concentration.

To bypass this barrier, we introduce a \textbf{3-Regime Structural Partition}:
Let $0 < 1/n_0 \ll \alpha \ll \beta \ll \delta \ll \epsilon \ll 1$. We decompose the tree $T$ into three vertex-induced regimes:
\begin{enumerate}
    \item \textbf{Super-Stars ($T_{\text{high}}$)}: Vertices with $\deg(v) \ge \Delta_{\text{high}} := \alpha n$. There are at most $1/\alpha = O(1)$ such hubs.
    \item \textbf{Intermediate Absorbers ($T_{\text{mid}}$)}: Vertices with $\Delta_{\text{mid}} \le \deg(v) < \Delta_{\text{high}}$, where $\Delta_{\text{mid}} := n / \log^2 n$.
    \item \textbf{Bounded Forest ($T_{\text{low}}$)}: The remaining forest with $\deg(v) < \Delta_{\text{mid}}$.
\end{enumerate}

\section{The Sequence of Seven Lemmas}

The logical skeleton of the proof rests on the following sequence of lemmas:

\begin{lemma}[Lemma 1: Reservoir and Absorber Construction]\label{lem:1}
There exists a buffer subset $R \subset \{0, \dots, \lfloor(1+\epsilon)(n-1)\rfloor\}$ of size $|R| \ge \epsilon n$ and an absorber subtree $T_{\text{abs}} \subset T$ of size $|V(T_{\text{abs}})| \le \delta n$ such that $R$ contains a family of length-7 switching paths capable of absorbing any deficit set $D_{\text{rem}}$ of size $\le \delta n$.
\end{lemma}

\begin{lemma}[Lemma 2A: High-Degree Deterministic Star Embedding]\label{lem:2a}
The components $T_{\text{high}}$ admit a deterministic injection into the extremal boundaries $\{0, 1, \dots\}$ and $\{\lfloor(1+\epsilon)(n-1)\rfloor, \dots\}$ using alternating $\alpha$-valuation coordinates, covering the largest difference band $[(1+\epsilon)n - |E(T_{\text{high}})|, (1+\epsilon)n]$ without collisions.
\end{lemma}

\begin{lemma}[Lemma 2B: Low-Degree LLL Embedding]\label{lem:2b}
For the forest $T_{\text{low}}$, because the local dependency degree satisfies $D \le 2\Delta_{\text{mid}} \le 2n/\log^2 n$, the Lov\'asz Local Lemma condition $e \cdot p \cdot D \le e (1/n)(2n/\log^2 n) \to 0$ holds. Hence, $T_{\text{low}}$ can be embedded quasi-randomly without edge difference collisions.
\end{lemma}

\begin{lemma}[Lemma 1-Cut: Interface Cut-Edge Separation]\label{lem:cut}
Let $E_{\text{cut}}$ be the edges joining $T_{\text{high}}$ and $T_{\text{low}}$. Because $|E_{\text{cut}}| \le O(\beta n) \ll \epsilon n$, there exist phase shifts $s_i \in R$ on boundary vertices such that all cut-edge differences fall disjointly into the intermediate window $I_{\text{mid}} := (|E(T_{\text{low}})|+\delta n, (1+\epsilon)n - |E(T_{\text{high}})|)$.
\end{lemma}

\begin{lemma}[Lemma 2-Parity: Bipartite Slack Balancing]\label{lem:parity}
Let $V(T) = A \cup B$ be the natural bipartition of $T$. The slack reservoir $R$ can be partitioned into $R_{\text{even}}$ and $R_{\text{odd}}$ such that $|R_{\text{even}}|/|R_{\text{odd}}| = |A|/|B| + O(1/\sqrt{n})$. Under this ratio, switching paths maintain parity neutrality, precluding parity deadlocks during absorption.
\end{lemma}

\begin{lemma}[Lemma 2C: Intermediate Absorption and Completion]\label{lem:2c}
By Hall's Marriage Theorem applied to the switching templates, the intermediate components $T_{\text{mid}}$ absorb all remaining missing differences $D_{\text{rem}}$, completing the edge set $E(T)$ with mutually distinct differences.
\end{lemma}

\begin{lemma}[Lemma 4: Slack Discarding and Injectivity]\label{lem:4}
The unused labels in $R \setminus R_{\text{used}}$ are discarded. Since $f$ is an injection into $\{0, \dots, \lfloor(1+\epsilon)(n-1)\rfloor\}$, the existence of unused labels introduces no edge conflicts.
\end{lemma}

\section{The $\epsilon \to 0$ Singularity}

When $\epsilon \to 0$, the slack reservoir collapses ($|R| \to 0$), and injectivity is replaced by rigid bijection. To bridge this gap, we conjecture the following algebraic closure:

\begin{conjecture}[Zero-Slack Boundary Inversion Conjecture]\label{conj:zeroslack}
For any finite defect set of edges $E_{\text{defect}}$ in a near-graceful tree, there exists a local 2-opt transposition chain of depth at most 2 on the neighborhood of leaves that cancels all difference collisions without requiring external buffer vertices.
\end{conjecture}

\section{Human-AI Collaborative Methodology}

This project was conducted as an open, cross-agent collaborative experiment:
\begin{itemize}
    \item \textbf{Human Director (kuji-nine)}: Provided problem selection, architectural trajectory, operational intuition, and cross-model synthesis.
    \item \textbf{Qwen3 827B (Local)}: Conducted structural inference, discrete search of minimal failure configurations, and counterexample discovery.
    \item \textbf{Gemini}: Formulated extremal graph-theoretic arguments, analyzed the DEM martingale breakdown, and authored the 7-lemma structural chain.
    \item \textbf{Claude (Whale Director)}: Orchestrated workflow execution, performed rigorous peer review and error correction, and authored formal documentation.
\end{itemize}

\section*{Availability and Archival}
The complete LaTeX sources, verification scripts, and raw execution logs are openly archived at:
\begin{center}
\url{https://github.com/kuji-nine/graceful-tree-theorem-strategy}
\end{center}
and citable via Zenodo DOI upon release.

\bibliographystyle{plain}
\begin{thebibliography}{99}
\bibitem{adamaszek2020almost}
M.~Adamaszek, P.~Allen, C.~Grosu, and J.~Hladk\'y.
\newblock Almost graceful labellings of trees.
\newblock {\em Random Structures \& Algorithms}, 56(3):637--668, 2020.

\bibitem{fang2010computational}
W.~Fang.
\newblock A computational approach to the graceful tree conjecture.
\newblock {\em arXiv preprint arXiv:1003.3045}, 2010.

\bibitem{montgomery2020proof}
R.~Montgomery, A.~Pokrovskiy, and B.~Sudakov.
\newblock A proof of {R}ingel's conjecture.
\newblock {\em Geometric and Functional Analysis}, 31(3):663--720, 2021.
\newblock (arXiv:2001.02665).

\bibitem{rosa1966certain}
A.~Rosa.
\newblock On certain valuations of the vertices of a graph.
\newblock {\em Theory of Graphs (Internat. Sympos., Rome)}, pages 349--355, 1966.
\end{thebibliography}

\end{document}
